\documentclass{article}
\usepackage[utf8]{inputenc}
\usepackage[T1]{fontenc}
\usepackage{geometry}
\usepackage{graphicx}
\usepackage{float}
\usepackage{hyperref}
\usepackage[english,indonesian]{babel}

\title{Optimasi Pemrosesan Template LaTeX Menggunakan Pendekatan Berbasis AI}
\author{Academic Support Agent}
\date{\today}

\begin{document}
\maketitle

\begin{abstract}

\end{abstract}

\chapter{PENDAHULUAN}


\chapter{TINJAUAN PUSTAKA}


% Original content (commented out):
% \documentclass[a4paper]{article}

\usepackage{fullpage} % Package to use full page
\usepackage{parskip} % Package to tweak paragraph skipping
\usepackage{tikz} % Package for drawing
\usepackage{amsmath}
\usepackage{hyperref}


\usepackage{biblatex}
\addbibresource{bibliography.bib} % Nodwch fod hwn yn wahanol i'r fersiwn Saesneg - er mwyn medu rhoi teitl Cymraeg i'r Cyferirnodau

\title{Deall y Deilliad}
\author{Vincent Knight}
\date{1984/02/14}

\renewcommand{\figurename}{Ffigwr} % Ail-enwi'r amgylchedd 'Figure' i 'Ffigwr'

\begin{document}

\maketitle

\section{Cyflwyniad}

Mae differu yn gysyniad yn Fathemateg a astudiwyd yng Nghalcwlws. Mae yna drafodaeth barhaus i weld pwy oedd y cyntaf i ddiffinio differu: Leibniz neu Newton \cite{bardi2006calculus}.

Mae differu yn galluogi cyfrifiad graddiant tangiad cromlin ar unrhyw bwynt, fel a gweler yn Ffigwr \ref{exampleplot}.


\begin{figure}[!htbp]
\begin{center}
\begin{tikzpicture}
\draw[domain=-2:2, color=blue] plot (\x, {1 - (\x)^2}) node[above = .5cm, right, color=blue] {$f(x)=1-x^2$};
\draw[domain=-2:2, color=red] plot(\x,-1 * \x + 1.25) node[above = .5cm, right, color=red] {Tangiad wrth $x=.5$};
\draw [thick, ->] (-3,0) -- (3,0) node [above] {$x$};
\draw [thick, ->] (0,-3) -- (0,3) node [right] {$y$};
\node at (.5,.75) {\textbullet};
\end{tikzpicture}
\end{center}
\caption{Y plot $f(x)=1-x^2$ gyda thangiad wrth $x=.5$.}\label{exampleplot}
\end{figure}

Mae differu nawr yn dechneg sy'n cael ei ddysgu i fyfyrwyr mathemateg trwy gydol y byd. Yn y ddogfen hon byddaf yn trafod rhai agweddau differu.


\section{Archwilio'r deilliad yn defnyddio Sage}

Diffiniad y terfan $f(x)$ yn $x=a$ wedi dynodi gan $f'(a)$ yw:

\begin{equation}
f'(a) = \lim_{h\to0}\frac{f(a+h)-f(a)}{h}
\end{equation}

Fe allwch ddefnyddio'r cod canlynol yn sage i roi'r terfan uchod:

\begin{verbatim}
def illustrate(f, a):
    """
    Ffwythiant sy'n cymryd ffwythiant ac yn dangos diffiniad terfannol y deilliad ar bwynt a rhoddir.
    """
    lst = []
    for h in srange(.01, 3, .01):
    	lst.append([h,(f(a+h)-f(a))/h])
    return list_plot(lst, axes_labels=['$x$','$\\frac{f(%.02f+h)-f(%.02f)}{h}$' % (a,a)])
\end{verbatim}

\begin{figure}[!htbp]
\begin{center}
\includegraphics[width=8cm]{sage1.png}
\end{center}
\caption{Y deilliad $f(x)=1-x^2$ yn $x=.5$ yn cydgyfeirio i -1 wrth i $h\to0$.}
\end{figure}

Os ydym ni eisiau plotio'r tangiad ym mhwynt $\alpha$ o'r ffwythiant, fe allwn ni defnyddio'r canlynol:

\begin{align}
y=&ax+b&&\text{(diffiniad llinell syth)}\nonumber\\
  &f'(a)x+b&&\text{(diffiniad y deilliad)}\nonumber\\
  &f'(a)x+f(a)-f'(a)a&&\text{(rydym yn gwybod fod y llinell yn croestorri $f$ yn $(a,f(a))$}\nonumber
\end{align}

Fe allwn gyfuno'r dull hwn gyda darn o god blaenorol i weld sut mae'r llinell tangiad yn cydgyfeirio wrth i diffiniad terfannol y deilliad cydgyfeirio:


\begin{verbatim}
def convergetangentialline(f, a, x1, x2, nbrofplots=50, epsilon=.1):
    """
    Ffwythiant sy'n gwneud i linell tangiad cydgyfeirio
    """
    clrs = rainbow(nbrofplots)
    k = 0
    h = epsilon
    p = plot(f, x, x1, x2)
    while k < nbrofplots:
        tangent(x) = fdash(f, a, h) * x + f(a) - fdash(f, a, h) * a
        p += plot(tangent(x), x, x1, x2, color=clrs[k])
        h += epsilon
        k += 1
    return p
\end{verbatim}

Mae'r plot a ddangosir yn Ffigwr \ref{lines} yn dangos sut mae'r llinellau yn cydgyfeirio i'r tangiad go iawn $1-x^2$ wrth i $x=2$ (y llinell goch yw'r gromlin agosaf).

\begin{figure}[!htbp]
\begin{center}
\includegraphics[width=8cm]{sage0.png}
\end{center}
\caption{Llinellau yn cydgyfeirio i'r tangiad wrth i $h\to0$.}\label{lines}
\end{figure}

Nodwch fod y plot olaf yn defnyddio diffiniad \textbf{go iawn} y deilliad ac nid y brasamcan.

\section{Casgliadau}

Yn yr adroddiad hwn rydw i wedi archwilio diffiniad y terfan, ac yn delweddu deilliad ffwythiant wrth i $h\to 0$. Mae'r cod a ddefnyddiwyd \url{https://sage.maths.cf.ac.uk/home/pub/18/} yn defnyddio gallu differu Sage yn ogystal a'i gallu plotio.

Mae yna agweddau amrywiol eraill gallaf wedi archwilio fel rheolau differu symbolaidd. Er enghraifft:

$$\frac{dx^n}{dx}=(n+1)x^{n}\text{ os yw }x\ne-1$$

Yn ogystal \^{a} hwn mae'n ddiddorol i nodi bodolaeth ffwythiannau sydd \textbf{ddim} yn ddifferadwy ym mhwynt, er enghraifft y ffwythiant $f(x)=\sin(1/x)$ sydd ddim yn ddifferadwy yn $x=0$. Dangosir plot o'r ffwythiant hwn yn Ffigwr \ref{notdiff}.

\begin{figure}[!htbp]
\begin{center}
\includegraphics[width=8cm]{sage2.png}
\end{center}
\caption{Ffwythiant annifferadwy yn $x=0$.}\label{notdiff}
\end{figure}


\printbibliography[title=Cyfeirnodau]
\end{document}

\end{document}